\documentclass[11pt]{article}
\usepackage[margin=1in]{geometry}
\usepackage{amsmath,amssymb}
\setlength{\parindent}{0pt}
\setlength{\parskip}{0.7em}
\begin{document}
\textbf{21-127 Concepts of Mathematics (fictional)}

Homework 1 - Foundations of proof --- Practice submission B

Student: Fictional learner

\section*{Question 1: Sets in three forms}

(a) $A=\{16,25,36,49,64,81\}$.

(b) $B=\{5k+1:k\in\mathbb{Z},\ k\geq0\}$.

(c) $C=\{3^k:k\in\mathbb{Z}\}$.

\newpage
\textbf{21-127 Concepts of Mathematics (fictional)}

Homework 1 - Foundations of proof --- Practice submission B

Student: Fictional learner

\section*{Question 2: A divisibility claim}

The statement is true. If a number divides a product, it must divide at least one factor. Applying that fact to $10$ gives the result.

\newpage
\textbf{21-127 Concepts of Mathematics (fictional)}

Homework 1 - Foundations of proof --- Practice submission B

Student: Fictional learner

\section*{Question 3: A rational point inside an interval}

The rational numbers are closed under addition, multiplication, and division by a nonzero rational. Therefore $u$ is rational.

Also $u-x=(y-x)/4>0$ and $y-u=3(y-x)/4>0$. Hence $x<u<y$.

\newpage
\textbf{21-127 Concepts of Mathematics (fictional)}

Homework 1 - Foundations of proof --- Practice submission B

Student: Fictional learner

\section*{Question 4: An odd expression}

For $n=0,1,2$, the expression is $1,5,11$. All three values are odd, and the outputs are increasing, so they continue to be odd for all integers $n$.

\newpage
\textbf{21-127 Concepts of Mathematics (fictional)}

Homework 1 - Foundations of proof --- Practice submission B

Student: Fictional learner

\section*{Question 5: Removing either set}

Take $t\in A\setminus(B\cup C)$. Then $t\in A$ and $t$ is in neither $B$ nor $C$. Hence $t\in A\setminus B$ and $t\in A\setminus C$. Therefore the two sets are equal.

\end{document}
